\documentclass[11pt]{article}

\usepackage[utf8]{inputenc}
\usepackage[T1]{fontenc}
\usepackage{lmodern}
\usepackage{geometry}
\usepackage{amsmath,amssymb,amsfonts,amsthm,mathtools}
\usepackage{bm}
\usepackage{enumitem}
\usepackage{microtype}
\usepackage{hyperref}
\usepackage{titlesec}
\usepackage{booktabs}
\usepackage{array}
\usepackage{setspace}
\usepackage{mathrsfs}
\usepackage{physics}

\geometry{margin=1in}
\hypersetup{colorlinks=true, linkcolor=black, urlcolor=blue, citecolor=black}
\setlist[enumerate]{topsep=0.4em,itemsep=0.35em}
\setlist[itemize]{topsep=0.3em,itemsep=0.25em}
\titleformat{\section}{\large\bfseries}{\thesection.}{0.5em}{}
\titleformat{\subsection}{\normalsize\bfseries}{\thesubsection.}{0.5em}{}
\setstretch{1.06}

\newcommand{\Aether}{\AE{}ther}
\newcommand{\Aflow}{\AE{}ther-flow}
\newcommand{\TheoryName}{\Aflow{} Interpretation of Relativity}
\newcommand{\TheoryProgram}{\Aether{} / \Aflow{} framework}

\newcommand{\CompletedMetric}{g^{\sharp}}
\newcommand{\MatterSpace}{\Psi}

\newcommand{\Car}{\mathrm{car}}
\newcommand{\Cur}{\mathrm{cur}}
\newcommand{\Hom}{\mathrm{hom}}

\newcommand{\ForSpace}[1]{\mathcal I_{#1}^{\mathrm{for}}}
\newcommand{\PrimShell}{\mathcal S_{\mathrm{prim}}}
\newcommand{\Reduction}[1]{\mathcal R_{#1}}
\newcommand{\CurrentCarrier}[1]{\mathfrak C_{#1}^{\mathrm{cur}}}
\newcommand{\EqCarrier}[1]{\mathfrak C_{#1}^{\mathrm{eq}}}
\newcommand{\MetricOut}{\mathscr G_{\mathrm{out}}}
\newcommand{\MatterOut}{\mathscr T_{\mathrm{out}}}

\newcommand{\SubTarget}[1]{E_{#1}^{\mathrm{sub}}}
\newcommand{\SubPair}[1]{\mathfrak s_{#1}^{\mathrm{sub}}}
\newcommand{\EqnTarget}[1]{Q_{#1}^{\mathrm{eqn}}}
\newcommand{\HiddenSpace}[1]{\mathcal H_{#1}^{\mathrm{hid}}}
\newcommand{\HiddenBody}[1]{D_{#1}^{\mathrm{hid}}}

\newcommand{\ProtoSpace}[1]{\mathcal U_{#1}^{\mathrm{proto}}}
\newcommand{\ProtoBody}[1]{\mathcal B_{#1}^{\mathrm{proto}}}
\newcommand{\ProtoSection}[1]{\Gamma_{#1}^{\mathrm{proto}}}
\newcommand{\ProtoShellSet}[1]{\mathcal Z_{#1}^{\mathrm{proto}}}
\newcommand{\ProtoZeroCore}[1]{\mathcal Z_{#1}^{0,\mathrm{proto}}}
\newcommand{\ProtoSplit}[1]{\Phi_{#1}^{\mathrm{proto}}}
\newcommand{\ProtoCharge}[1]{\mathcal Q_{#1}^{\mathrm{proto}}}
\newcommand{\ProtoEmbed}[1]{\zeta_{#1}^{\mathrm{proto}}}
\newcommand{\ProtoKernelCh}[1]{\widetilde K_{#1}^{0,\mathrm{proto}}}
\newcommand{\ProtoStateComp}[1]{P_{#1}^{\mathrm{proto,co}}}

\newcommand{\PreProtoNucleus}{\mathfrak Q^{\mathrm{nuc}}}
\newcommand{\PreProtoSpace}[1]{\mathcal V_{#1}^{\mathrm{preproto}}}
\newcommand{\PreProtoBody}[1]{\mathcal B_{#1}^{\mathrm{preproto}}}
\newcommand{\PreProtoProj}[1]{\Pi_{#1}^{\mathrm{preproto}}}
\newcommand{\PreProtoProtoMin}[1]{\mathfrak f_{#1}^{\mathrm{preproto,proto}}}
\newcommand{\PreProtoSection}[1]{\Gamma_{#1}^{\mathrm{preproto}}}
\newcommand{\PreProtoFunctional}[1]{\mathscr W_{#1}^{\mathrm{preproto}}}
\newcommand{\PreProtoShellSet}[1]{\mathcal Z_{#1}^{\mathrm{preproto}}}

\newcommand{\PreNucPackage}{\mathfrak R^{\mathrm{prenuc}}}
\newcommand{\PreNucCore}{\mathfrak R^{\mathrm{core}}}
\newcommand{\PreNucSpace}[1]{\mathcal W_{#1}^{\mathrm{prenuc}}}
\newcommand{\PreNucBody}[1]{\mathcal B_{#1}^{\mathrm{prenuc}}}
\newcommand{\PreNucProj}[1]{\Pi_{#1}^{\mathrm{prenuc}}}
\newcommand{\PreNucNucleusMin}[1]{\mathfrak f_{#1}^{\mathrm{prenuc,nuc}}}
\newcommand{\PreNucSection}[1]{\Gamma_{#1}^{\mathrm{prenuc}}}
\newcommand{\PreNucFunctional}[1]{\mathscr W_{#1}^{\mathrm{prenuc}}}
\newcommand{\PreNucShellSet}[1]{\mathcal Z_{#1}^{\mathrm{prenuc}}}

\newcommand{\PreCorePackage}{\mathfrak S^{\mathrm{precore}}}
\newcommand{\PreCoreSpace}[1]{\mathcal Y_{#1}^{\mathrm{precore}}}
\newcommand{\PreCoreBody}[1]{\mathcal B_{#1}^{\mathrm{precore}}}
\newcommand{\PreCoreProj}[1]{\Pi_{#1}^{\mathrm{precore}}}
\newcommand{\PreCoreCoreMin}[1]{\mathfrak f_{#1}^{\mathrm{precore,core}}}
\newcommand{\PreCoreShellSet}[1]{\mathcal Z_{#1}^{\mathrm{precore}}}

\newtheorem{definition}{Definition}
\newtheorem{proposition}{Proposition}
\newtheorem{remark}{Remark}
\newtheorem{corollary}{Corollary}
\newtheorem{theorem}{Theorem}

\title{\textbf{The \TheoryName{}}\\[0.4em]
\large Primitive-Reservoir Observer-Reduction\\
\large Observer-Relevant Pre-Nucleus Exact-Shell Core\\
\large from Exact-Shell-Core-Generating Substrate Pre-Core Dynamics}
\author{}
\date{}

\begin{document}

\maketitle

\begin{abstract}
The current primitive-reservoir observer-reduction frontier already spent the honest internal collapse available at full pre-nucleus scope. The active exact-shell-core theorem showed that the full nucleus-generating substrate pre-nucleus package contains benchmark-neutral ambient presentation data and that the live burden factors instead through a smaller observer-relevant exact-shell core. After that step, another same-output deeper-origin relay that merely restates that same exact-shell core is no longer the honest default move. The next admissible benchmark-facing gain must either derive or justify that core from still more primitive substrate structure, or identify a genuinely smaller theorem-level observer-relevant datum strictly inside it.

The present note takes the first route. For each sector it introduces \emph{exact-shell-core-generating substrate pre-core dynamics}: a finite-dimensional pre-core state space, a compact convex pre-core body over the recorded exact-shell-core body, a smooth pre-core minimizer on that fixed body, and a closed exact pre-core shell projecting diffeomorphically to the recorded exact-shell-core shell. The current observer-relevant pre-nucleus exact-shell core is then recovered canonically as the projected exact-shell transport image of that deeper pre-core package.

The benchmark boundary remains fixed throughout. The one operative metric is still exactly \(\CompletedMetric\), the ordinary matter route is unchanged, there is no second operative metric, the low-energy matter law is not enlarged, and no stronger promotion language is licensed. The positive result is therefore narrow but benchmark facing: the current exact-shell core is no longer left as the primitive stopping point on the active line. The remaining honest burden moves one layer deeper again, to the exact-shell-core-generating substrate pre-core dynamics themselves or to a sharper collapse strictly inside that package.
\end{abstract}

\section{Introduction}

The active derivational program of \TheoryName{} has again reached a stage where the honest next move is constrained more sharply than ``go one layer deeper.'' The current active-primary theorem already removed the full pre-nucleus package as the live stopping point by proving that only a smaller observer-relevant exact-shell core remains benchmark-facingly live there. That result also fixed the admissible next alternatives: derive or justify that exact-shell core from still more primitive substrate structure, or prove a genuinely smaller collapse strictly inside it.

The present note takes the default first route. Its role is not to reopen the full pre-nucleus package whose discarded ambient presentation data have already been screened as benchmark neutral. Its role is to show that even the now-live exact-shell core can itself be presented as the projected exact-shell transport image of a deeper substrate package. In that sense the current exact-shell core is no longer the primitive place where the active line has to stop.

Two features matter here. First, the theorem targets the exact-shell core itself rather than returning to the already-screened larger pre-nucleus package, so it respects the routing safeguard against recursive depth drift. Second, the deeper pre-core package is not merely another renaming of the current core. It adds an explicit pre-core state space, pre-core body, pre-core-to-core projection, pre-core minimizer, and deeper exact-shell transport below the current frontier. The current exact-shell core is therefore recovered from more primitive substrate data rather than simply reasserted.

\paragraph{Framework and claim boundary.}
\TheoryProgram{} denotes the broader \Aether{} / \Aflow{} framework built on the claims that the \Aether{} is the underlying four-dimensional substrate of reality, the \Aflow{} is its intrinsic ordered motion, observed three-dimensional space is the local experiential slice of that deeper substrate, S-time is the experienced order of change arising from matter, light, and the \Aflow{}, and observed expansion is the three-dimensional appearance of deeper four-dimensional ordered motion. Within that framework, \TheoryName{} denotes the exact-closure theory in which the effective gravitational dynamics are adopted to be exactly Einsteinian with universal matter coupling. In the active sequence, \emph{adoption} means use of that established relativistic dynamics without claiming substrate derivation, while \emph{derivation} is reserved for a first-principles recovery from explicit substrate variables.

The present note stays strictly inside that boundary. It does not alter the benchmark exact-closure package. It does not introduce a second operative metric or a new low-energy matter law. It only proves that the current observer-relevant pre-nucleus exact-shell core can itself be generated from a more primitive pre-core substrate dynamics.

\section{Fixed Inputs}

\begin{definition}[Recorded primitive continuation data]
\label{def:fixed}
Fix the following continuation data from the active primitive-reservoir observer-reduction line.
\begin{enumerate}
    \item For each sector label \(\sigma\in\{\Car,\Cur,\Hom\}\), a primitive forcing shell
    \[
    \ForSpace{\sigma},
    \]
    a visible reduction map
    \[
    \Reduction{\sigma}:\ForSpace{\sigma}\to X_{\sigma},
    \]
    and shell-level current and equilibrium carriers
    \[
    \CurrentCarrier{\sigma}:\ForSpace{\sigma}\to Z_{\sigma}^{\mathrm{cur}},
    \qquad
    \EqCarrier{\sigma}:\ForSpace{\sigma}\to Z_{\sigma}^{\mathrm{eq}}.
    \]
    \item The product shell
    \[
    \PrimShell
    \equiv
    \ForSpace{\Car}\times \ForSpace{\Cur}\times \ForSpace{\Hom}.
    \]
    \item The recorded metric and ordinary-matter outputs
    \[
    \MetricOut:\PrimShell\to\{\text{observer metrics}\},
    \qquad
    \MatterOut:\PrimShell\times\MatterSpace\to\{\text{observer stress tensors}\},
    \]
    satisfying
    \begin{equation}
    \MetricOut(\mathbf w)=\CompletedMetric,
    \qquad
    \MatterOut(\mathbf w,\psi)
    =
    T_{\mu\nu}^{\mathrm{matter}}[\CompletedMetric,\psi]
    \label{eq:fixedoutputs}
    \end{equation}
    for every \(\mathbf w\in\PrimShell\) and every matter configuration \(\psi\in\MatterSpace\).
\end{enumerate}
\end{definition}

\begin{definition}[Recorded proto-germ exact-shell target]
\label{def:proto}
Fix \(\sigma\in\{\Car,\Cur,\Hom\}\). The active line already fixes:
\begin{enumerate}
    \item the same substrate target and pairing
    \[
    \SubTarget{\sigma},
    \qquad
    \SubPair{\sigma}:
    \SubTarget{\sigma}\times\SubTarget{\sigma}\to\mathbb R;
    \]
    \item a hidden fiber space \(\HiddenSpace{\sigma}\) with compact hidden body \(\HiddenBody{\sigma}\subset\HiddenSpace{\sigma}\);
    \item a finite-dimensional proto-germ state space \(\ProtoSpace{\sigma}\), a compact convex proto-germ body \(\ProtoBody{\sigma}\subset\ProtoSpace{\sigma}\), a smooth proto-germ restriction section
    \[
    \ProtoSection{\sigma}:\ProtoSpace{\sigma}\to\SubTarget{\sigma},
    \]
    and the exact proto-germ shell
    \[
    \ProtoShellSet{\sigma}
    \equiv
    \bigl(\ProtoSection{\sigma}\bigr)^{-1}(0)\cap\ProtoBody{\sigma};
    \]
    \item a closed embedded proto zero core \(\ProtoZeroCore{\sigma}\subset\ProtoShellSet{\sigma}\) and a split diffeomorphism
    \[
    \ProtoSplit{\sigma}:
    \ProtoZeroCore{\sigma}\times\HiddenBody{\sigma}
    \xrightarrow{\cong}
    \ProtoShellSet{\sigma};
    \]
    \item a hidden charge
    \[
    \ProtoCharge{\sigma}:\HiddenSpace{\sigma}\to\EqnTarget{\sigma},
    \]
    a smooth observer-compatible exact proto-germ zero-response realization
    \[
    \ProtoEmbed{\sigma}:\ForSpace{\sigma}\hookrightarrow\ProtoZeroCore{\sigma},
    \]
    and a fixed-data solvability / coercive package on \(\ProtoZeroCore{\sigma}\), recorded through \(\ProtoKernelCh{\sigma}\) and \(\ProtoStateComp{\sigma}\).
\end{enumerate}
\end{definition}

\begin{definition}[Recorded current transport nucleus]
\label{def:nucleus}
Fix \(\sigma\in\{\Car,\Cur,\Hom\}\). The recorded current transport nucleus \(\PreProtoNucleus\) for sector \(\sigma\) consists of:
\begin{enumerate}
    \item the pre-proto-germ state space \(\PreProtoSpace{\sigma}\), compact convex body \(\PreProtoBody{\sigma}\subset\PreProtoSpace{\sigma}\), projection
    \[
    \PreProtoProj{\sigma}:\PreProtoSpace{\sigma}\to\ProtoSpace{\sigma},
    \]
    and proto section
    \[
    \PreProtoProtoMin{\sigma}:\ProtoBody{\sigma}\to\PreProtoBody{\sigma}
    \]
    satisfying
    \[
    \PreProtoProj{\sigma}\bigl(\PreProtoBody{\sigma}\bigr)=\ProtoBody{\sigma},
    \qquad
    \PreProtoProj{\sigma}\circ\PreProtoProtoMin{\sigma}
    =
    \mathrm{id}_{\ProtoBody{\sigma}};
    \]
    \item the restriction section
    \[
    \PreProtoSection{\sigma}:\PreProtoSpace{\sigma}\to\SubTarget{\sigma},
    \]
    the induced functional
    \[
    \PreProtoFunctional{\sigma}(q)
    \equiv
    \SubPair{\sigma}\bigl(\PreProtoSection{\sigma}(q),\PreProtoSection{\sigma}(q)\bigr),
    \]
    and the exact shell
    \[
    \PreProtoShellSet{\sigma}
    \equiv
    \bigl(\PreProtoSection{\sigma}\bigr)^{-1}(0)\cap\PreProtoBody{\sigma},
    \]
    with
    \[
    \PreProtoSection{\sigma}\circ\PreProtoProtoMin{\sigma}
    =
    \ProtoSection{\sigma},
    \qquad
    \PreProtoProj{\sigma}\big|_{\PreProtoShellSet{\sigma}}
    :
    \PreProtoShellSet{\sigma}
    \xrightarrow{\cong}
    \ProtoShellSet{\sigma};
    \]
    \item benchmark faithfulness, meaning preservation of the benchmark outputs \eqref{eq:fixedoutputs}, no second operative metric, no enlarged low-energy matter law, and no premature promotion from adoption to completed derivation.
\end{enumerate}
\end{definition}

\begin{definition}[Recorded observer-relevant pre-nucleus exact-shell core]
\label{def:core}
Fix \(\sigma\in\{\Car,\Cur,\Hom\}\). The \emph{observer-relevant pre-nucleus exact-shell core} \(\PreNucCore\) for sector \(\sigma\) consists of:
\begin{enumerate}
    \item the pre-nucleus state space \(\PreNucSpace{\sigma}\), compact convex body \(\PreNucBody{\sigma}\subset\PreNucSpace{\sigma}\), affine surjection
    \[
    \PreNucProj{\sigma}:\PreNucSpace{\sigma}\to\PreProtoSpace{\sigma},
    \]
    and fiber minimizer
    \[
    \PreNucNucleusMin{\sigma}:\PreProtoBody{\sigma}\to\PreNucBody{\sigma}
    \]
    satisfying
    \[
    \PreNucProj{\sigma}\bigl(\PreNucBody{\sigma}\bigr)=\PreProtoBody{\sigma},
    \qquad
    \PreNucProj{\sigma}\circ\PreNucNucleusMin{\sigma}
    =
    \mathrm{id}_{\PreProtoBody{\sigma}};
    \]
    \item a closed embedded exact pre-nucleus shell
    \[
    \PreNucShellSet{\sigma}\subset\PreNucBody{\sigma}
    \]
    satisfying
    \[
    \PreNucNucleusMin{\sigma}\bigl(\PreProtoShellSet{\sigma}\bigr)
    \subset
    \PreNucShellSet{\sigma},
    \qquad
    \PreNucProj{\sigma}\big|_{\PreNucShellSet{\sigma}}
    :
    \PreNucShellSet{\sigma}
    \xrightarrow{\cong}
    \PreProtoShellSet{\sigma};
    \]
    \item benchmark faithfulness in the sense of Definition \ref{def:nucleus}.
\end{enumerate}
The ambient pre-nucleus restriction section \(\PreNucSection{\sigma}\) and the induced quadratic functional \(\PreNucFunctional{\sigma}\) are not taken as independent core data.
\end{definition}

\begin{remark}
Definitions \ref{def:fixed}--\ref{def:core} are fixed inputs. The one operative metric, the ordinary matter route, the fixed proto-germ exact-shell target, the current transport nucleus, and the current observer-relevant exact-shell core are not reopened here. The present note asks only whether that current exact-shell core can itself be generated from more primitive substrate dynamics.
\end{remark}

\section{Exact-Shell-Core-Generating Substrate Pre-Core Dynamics}

\begin{definition}[Exact-shell-core-generating substrate pre-core dynamics]
\label{def:precore}
Fix \(\sigma\in\{\Car,\Cur,\Hom\}\). An \emph{exact-shell-core-generating substrate pre-core dynamics package} \(\PreCorePackage\) for sector \(\sigma\) consists of the following data.
\begin{enumerate}
    \item A finite-dimensional Euclidean pre-core state space \(\PreCoreSpace{\sigma}\), a compact convex pre-core body
    \[
    \PreCoreBody{\sigma}\subset\PreCoreSpace{\sigma}
    \]
    with nonempty interior, an affine surjection
    \[
    \PreCoreProj{\sigma}:\PreCoreSpace{\sigma}\to\PreNucSpace{\sigma},
    \]
    and a smooth pre-core minimizer
    \[
    \PreCoreCoreMin{\sigma}:\PreNucBody{\sigma}\to\PreCoreBody{\sigma}
    \]
    satisfying
    \[
    \PreCoreProj{\sigma}\bigl(\PreCoreBody{\sigma}\bigr)=\PreNucBody{\sigma},
    \qquad
    \PreCoreProj{\sigma}\circ\PreCoreCoreMin{\sigma}
    =
    \mathrm{id}_{\PreNucBody{\sigma}}.
    \]
    \item A closed embedded exact pre-core shell
    \[
    \PreCoreShellSet{\sigma}\subset\PreCoreBody{\sigma}
    \]
    satisfying
    \[
    \PreCoreCoreMin{\sigma}\bigl(\PreNucShellSet{\sigma}\bigr)
    \subset
    \PreCoreShellSet{\sigma},
    \qquad
    \PreCoreProj{\sigma}\big|_{\PreCoreShellSet{\sigma}}
    :
    \PreCoreShellSet{\sigma}
    \xrightarrow{\cong}
    \PreNucShellSet{\sigma}.
    \]
    \item Benchmark faithfulness, meaning preservation of the benchmark outputs \eqref{eq:fixedoutputs}, no second operative metric, no enlarged low-energy matter law, and no premature promotion from adoption to completed derivation.
\end{enumerate}
\end{definition}

\begin{remark}
Definition \ref{def:precore} is intentionally written over the recorded exact-shell-core target of Definition \ref{def:core}. It does not reintroduce the previously screened ambient pre-nucleus section or quadratic functional as live data. Its positive role is narrower: provide a more primitive substrate package whose projected body transport and exact-shell transport already generate the current exact-shell core.
\end{remark}

\section{Canonical Recovery of the Current Exact-Shell Core}

\begin{proposition}[Pre-core dynamics generate the recorded exact-shell core]
\label{prop:core}
Every benchmark-faithful pre-core package in the sense of Definition \ref{def:precore} generates the recorded observer-relevant pre-nucleus exact-shell core of Definition \ref{def:core} by projected body transport and exact-shell transport.
\end{proposition}

\begin{proof}
Definition \ref{def:precore}(1) gives a compact convex pre-core body together with an affine surjection
\[
\PreCoreProj{\sigma}:\PreCoreSpace{\sigma}\to\PreNucSpace{\sigma}
\]
and a smooth minimizer
\[
\PreCoreCoreMin{\sigma}:\PreNucBody{\sigma}\to\PreCoreBody{\sigma}
\]
satisfying
\[
\PreCoreProj{\sigma}\bigl(\PreCoreBody{\sigma}\bigr)=\PreNucBody{\sigma},
\qquad
\PreCoreProj{\sigma}\circ\PreCoreCoreMin{\sigma}
=
\mathrm{id}_{\PreNucBody{\sigma}}.
\]
Definition \ref{def:precore}(2) gives a closed embedded exact pre-core shell projecting diffeomorphically to the recorded exact-shell-core shell:
\[
\PreCoreProj{\sigma}\big|_{\PreCoreShellSet{\sigma}}
:
\PreCoreShellSet{\sigma}
\xrightarrow{\cong}
\PreNucShellSet{\sigma}.
\]
Therefore the recorded body transport and exact-shell transport of Definition \ref{def:core} arise as the projected exact-shell image of the deeper pre-core package. Benchmark faithfulness is preserved by Definition \ref{def:precore}(3).
\end{proof}

\begin{proposition}[Downstream transport to the recorded transport nucleus persists]
\label{prop:downstream}
Fix a benchmark-faithful pre-core package \(\PreCorePackage\). Then
\[
(\PreNucProj{\sigma}\circ\PreCoreProj{\sigma})\bigl(\PreCoreBody{\sigma}\bigr)
=
\PreProtoBody{\sigma},
\]
\[
(\PreNucProj{\sigma}\circ\PreCoreProj{\sigma})\circ
\PreCoreCoreMin{\sigma}\circ\PreNucNucleusMin{\sigma}
=
\mathrm{id}_{\PreProtoBody{\sigma}},
\]
and the composite shell transport
\[
(\PreNucProj{\sigma}\circ\PreCoreProj{\sigma})\big|_{\PreCoreShellSet{\sigma}}
:
\PreCoreShellSet{\sigma}
\xrightarrow{\cong}
\PreProtoShellSet{\sigma}
\]
is a diffeomorphism. Consequently
\[
(\PreProtoProj{\sigma}\circ\PreNucProj{\sigma}\circ\PreCoreProj{\sigma})\big|_{\PreCoreShellSet{\sigma}}
:
\PreCoreShellSet{\sigma}
\xrightarrow{\cong}
\ProtoShellSet{\sigma}
\]
is also a diffeomorphism.
\end{proposition}

\begin{proof}
Definition \ref{def:precore}(1) gives
\[
\PreCoreProj{\sigma}\bigl(\PreCoreBody{\sigma}\bigr)=\PreNucBody{\sigma}
\]
and
\[
\PreCoreProj{\sigma}\circ\PreCoreCoreMin{\sigma}
=
\mathrm{id}_{\PreNucBody{\sigma}}.
\]
Composing with the recorded exact-shell-core data of Definition \ref{def:core}(1) yields
\[
(\PreNucProj{\sigma}\circ\PreCoreProj{\sigma})\bigl(\PreCoreBody{\sigma}\bigr)
=
\PreNucProj{\sigma}\bigl(\PreNucBody{\sigma}\bigr)
=
\PreProtoBody{\sigma}
\]
and
\[
(\PreNucProj{\sigma}\circ\PreCoreProj{\sigma})\circ
\PreCoreCoreMin{\sigma}\circ\PreNucNucleusMin{\sigma}
=
\PreNucProj{\sigma}\circ\PreNucNucleusMin{\sigma}
=
\mathrm{id}_{\PreProtoBody{\sigma}}.
\]

For shell transport, Definition \ref{def:precore}(2) gives a diffeomorphism
\[
\PreCoreProj{\sigma}\big|_{\PreCoreShellSet{\sigma}}
:
\PreCoreShellSet{\sigma}
\xrightarrow{\cong}
\PreNucShellSet{\sigma},
\]
while Definition \ref{def:core}(2) gives a diffeomorphism
\[
\PreNucProj{\sigma}\big|_{\PreNucShellSet{\sigma}}
:
\PreNucShellSet{\sigma}
\xrightarrow{\cong}
\PreProtoShellSet{\sigma}.
\]
Their composition is therefore a diffeomorphism from \(\PreCoreShellSet{\sigma}\) to \(\PreProtoShellSet{\sigma}\). Composing once more with the shell diffeomorphism in Definition \ref{def:nucleus}(2) yields the final displayed diffeomorphism to \(\ProtoShellSet{\sigma}\).
\end{proof}

\begin{proposition}[The benchmark package remains fixed]
\label{prop:benchmark}
Every benchmark-faithful pre-core package preserves the same benchmark one-metric package \eqref{eq:fixedoutputs}. In particular, the operative metric remains exactly \(\CompletedMetric\), the ordinary matter route is unchanged, there is no second operative metric, and the low-energy matter law is not enlarged.
\end{proposition}

\begin{proof}
This is part of benchmark faithfulness in Definition \ref{def:precore}(3). The present note introduces only a deeper substrate-side source for the current exact-shell core. It does not alter the benchmark exact-closure package.
\end{proof}

\section{Main Theorem}

\begin{theorem}[Observer-relevant pre-nucleus exact-shell core from exact-shell-core-generating substrate pre-core dynamics]
\label{thm:main}
Fix the primitive continuation data of Definition \ref{def:fixed}, the recorded proto-germ exact-shell target of Definition \ref{def:proto}, the recorded current transport nucleus of Definition \ref{def:nucleus}, the recorded observer-relevant pre-nucleus exact-shell core of Definition \ref{def:core}, and a benchmark-faithful pre-core package in the sense of Definition \ref{def:precore}. Then, for each sector \(\sigma\in\{\Car,\Cur,\Hom\}\), the current exact-shell core is no longer an unexplained primitive stopping point on the active line: it is recovered canonically as the projected exact-shell transport image of the deeper pre-core package. More precisely:
\begin{enumerate}
    \item the pre-core package generates the recorded exact-shell core by Proposition \ref{prop:core};
    \item the exact-shell transport from the pre-core package already carries the recorded downstream transport to the current transport nucleus and the fixed proto-germ exact-shell target by Proposition \ref{prop:downstream};
    \item the benchmark boundary remains fixed by Proposition \ref{prop:benchmark};
    \item consequently the remaining honest burden below the previous exact-shell-core frontier is no longer the observer-relevant pre-nucleus exact-shell core \(\PreNucCore\) as such. What remains is to derive or justify the deeper exact-shell-core-generating substrate pre-core dynamics \(\PreCorePackage\) from still more primitive substrate structure, or else prove a still sharper collapse strictly inside that deeper package.
\end{enumerate}
\end{theorem}

\begin{proof}
Item (1) is Proposition \ref{prop:core}. Item (2) is Proposition \ref{prop:downstream}. Item (3) is Proposition \ref{prop:benchmark}. Item (4) is the routing consequence of the previous items together with the active safeguard against counting another same-output deeper-origin relay as new front-line progress when the live benchmark-relevant datum is unchanged.
\end{proof}

\begin{corollary}[What must be done next]
\label{cor:next}
After Theorem \ref{thm:main}, the next honest benchmark-facing manuscript must either:
\begin{enumerate}
    \item derive or justify the exact-shell-core-generating substrate pre-core dynamics \(\PreCorePackage\) from still more primitive substrate structure in a way that changes benchmark-facing status;
    \item identify a genuinely smaller theorem-level observer-relevant datum strictly inside \(\PreCorePackage\) and prove a sharper collapse to it;
    \item do either of the previous two while preserving the same benchmark one-metric package and the same claim boundary.
\end{enumerate}
Another same-output deeper-origin relay that preserves the same observer-relevant exact-shell core, another re-expansion back to the full pre-nucleus package as if its screened ambient presentation data were still live, another restatement of the pre-core package without a new derivational gain, another reopening of the former transport nucleus or larger pre-proto-germ packages as if they were still the frontier, or any move that introduces a second operative metric, enlarges the ordinary matter law, or uses stronger promotion language does not satisfy the present burden. If a quantum continuation is pursued, it matters benchmark-facingly only insofar as it derives this same exact-shell core, collapses the deeper pre-core package further, or closes a benchmark derivation gate in a genuinely new way.
\end{corollary}

\begin{proof}
Theorem \ref{thm:main} shows that the current exact-shell core is now a recovered target rather than the primitive stopping point of the active line. Therefore a subsequent manuscript must improve on that status rather than merely accompany it. What remains open is either to derive or justify the deeper pre-core package itself from more primitive structure, or to discover a smaller theorem-level observer-relevant datum inside it.
\end{proof}

\section{Conclusion}

The positive result of this note is a disciplined deeper derivation step below the current exact-shell-core frontier. The previous active-primary theorem already removed the full pre-nucleus package as the live burden by collapsing it to a smaller observer-relevant exact-shell core. The present theorem then takes the next honest route available there: it shows that even that current exact-shell core is not the primitive place where the active line has to stop, because it is recovered canonically from a deeper exact-shell-core-generating substrate pre-core dynamics.

The scope remains disciplined. The benchmark package is unchanged: the one operative metric is still exactly \(\CompletedMetric\), the ordinary matter route is unchanged, there is no second operative metric, and the low-energy matter law is not enlarged. Nothing here upgrades adoption to completed derivation or licenses stronger promotion language.

The open burden therefore moves one honest step further again. The next benchmark-facing manuscript should derive or justify the deeper pre-core package from still more primitive substrate structure, unless the sharper honest gain available instead is a still smaller theorem-level collapse strictly inside that package. Until then, the present result should be read for what it is: a benchmark-facing derivation of the current exact-shell core from deeper substrate dynamics within \TheoryName{}, not a claim that the full first-principles program is finished.

\end{document}
